\section{Targeted Practice Questions}\label{prac:questions}\label{prac:main}
These practice variations combine the supplied lectures, Tutorials 1--4,
and question styles from the two sample finals.  They are speculative
revision questions, not information about the actual midterm.  Marks are
suggested practice weights.  Attempt the questions before consulting the
separate answer sketches.

\subsection{Questions}
\paragraph{Q1. Scheduling with two levels of quotas (10 marks).}
There are workers $i$, disjoint periods $j$, and tasks, each belonging to
one period.  Worker $i$ can perform tasks in $A_i$, may perform at most
$b_i$ tasks overall, and at most one task per period.  Every task must be
performed by exactly one worker.
(a) Construct a flow network and specify the target flow value.
(b) Prove both directions of the reduction.  (c) Give an ILP.
(d) Explain why a fractional feasible flow is not itself a schedule, and
why this does not prevent a flow-based algorithm.

\paragraph{Q2. Bounded-frequency weighted covering (8 marks).}
Each element belongs to at most $k$ candidate sets with nonnegative costs.
Give an ILP and a deterministic $k$-approximation using its LP relaxation.
Prove feasibility and cost separately.  Does the same threshold proof work
if instead every \emph{set contains at most $k$ elements}?

\paragraph{Q3. Randomized cut, then remove the randomness (8 marks).}
For weighted MAX-$k$-CUT on a loopless undirected graph with nonnegative
weights, design an algorithm obtaining at least $(1-1/k)\OPT$ in expectation.
Convert it to a deterministic polynomial-time algorithm.  State precisely
how to compute conditional expectation after a partial labelling.  Explain
why independence of the edge indicators is unnecessary.

\paragraph{Q4. Randomized set cover with a stopping rule (10 marks).}
Solve the natural weighted-set-cover LP on a feasible instance with $n\ge1$
elements.  One trial takes $r=\lceil\ln(2n)\rceil$ independent rounds,
selecting each set $j$ independently with probability $x_j^*$ in each round
and taking their union.  Discard unsuccessful trials and repeat independently
until a cover is obtained.  Prove almost-sure feasibility, expected polynomial
running time, and expected returned cost at most $2r\OPT$.  Explain why
``one trial has expected cost at most $rLP^*$, so the returned cover does
too'' is not justified.

\paragraph{Q5. Minimum-cut structure and tie-breaking (8 marks).}
(a) Fix any maximum flow $f$.  Prove that for \emph{every} minimum cut,
outgoing arcs are saturated and incoming arcs carry zero flow.
(b) Deduce that a positive-capacity arc cannot cross two minimum cuts in
opposite directions.  (c) For integer capacities, find a minimum cut with
the fewest outgoing arcs.

\paragraph{Q6. MAX-DICUT rounding and a changed probability (8 marks).}
For a loopless directed graph with nonnegative arc weights, use the LP
$z_{ij}\le x_i$, $z_{ij}\le1-x_j$, with variables in $[0,1]$, maximizing
$\sum_{ij}w_{ij}z_{ij}$.  Prove that independently placing vertex $i$ in
the source side with probability $p_i=1/4+x_i/2$ gives at least half the
LP value in expectation.  As an extension, determine whether direct
rounding $p_i=x_i$ also guarantees half the optimum when the LP solution
is optimal.

\paragraph{Q7. MAX-SAT clause lengths and combination (10 marks).}
Write the weighted MAX-SAT ILP and LP.  Prove the LP-rounding bound for a
clause of length $\ell$, including the AM--GM and concavity steps.  Then
show why taking the better of LP rounding and fair random assignment
obtains $3\OPT/4$ in expectation.  State the preprocessing needed for
the clause-product calculation.

\paragraph{Q8. Metric TSP and the role of assumptions (8 marks).}
Prove the double-tree factor $2$ and Christofides factor $3/2$.  Explain
exactly where triangle inequality, even parity, and minimum-weight perfect
matching are used.  Give a four-vertex metric instance where a valid
double-tree execution is nonoptimal.  Explain why arbitrary nonmetric
costs invalidate the shortcut argument.

\paragraph{Q9. Knapsack scaling (8 marks).}
For nonnegative integer item values, derive the value-indexed $O(n^2V)$ DP,
where $V$ is the largest item value.  Use it to obtain an FPTAS with
$O(n^3/\epsilon)$ arithmetic-operation complexity.  Prove the guarantee
and explain why overweight items must be removed before defining $V$.

\paragraph{Q10. Vertex-disjoint routes (8 marks).}
For distinct nonadjacent vertices $x,y$ of an undirected graph, prove that
the maximum number of internally vertex-disjoint $x$--$y$ paths equals
the minimum size of an $x$--$y$ vertex separator.  Give the flow construction
and justify both correspondences.  Why do capacity-$1$ splitting arcs
suffice, and why are the other arcs assigned capacity larger than $n$?

\paragraph{Q11. Repair a flow algorithm (8 marks).}
(a) A scaling algorithm starts at $k=C$, halves by real division, and stops
below $1$.  Find an integer-capacity counterexample and repair the algorithm.
(b) Prove that BFS augmentations give $O(nm^2)$ time independent of capacity
magnitudes.  (c) Explain why arbitrary Ford--Fulkerson is only
pseudo-polynomial for integer capacities and may fail to terminate for
irrational capacities.

\paragraph{Q12. An unfamiliar objective with a familiar proof (6 marks).}
You are given $n$ strings of length $L$ over an alphabet of $q$ symbols.
Give a deterministic $1/q$ retained-value approximation for their longest
common subsequence in $O(nL+q)$ time.  Give a second proof using a sum of
per-symbol upper bounds on $\OPT$.

\subsection{Quick true/false drill}
For each statement give a proof or a small counterexample.
\begin{enumerate}[label=\textbf{T\arabic*.}]
  \item Complementing any 2-approximate vertex cover gives a constant-factor
        approximate maximum independent set.
  \item Integer capacities imply every maximum flow is integral on every arc.
  \item A feasible flow whose value equals the capacity of some cut is maximum.
  \item An integer-capacity flow algorithm taking $O(mC)$ time is polynomial
        in the binary input length.
  \item A constant-factor approximation guarantee for MAX-SAT automatically
        gives one for minimizing unsatisfied clauses.
  \item Shortest residual distances never decrease under arbitrary
        Ford--Fulkerson path choices.
  \item A feasible solution of a minimization LP relaxation upper-bounds
        the integer optimum.
  \item Every graph on $n\ge3$ vertices has vertex-cover LP integrality gap
        at least $2-2/n$.
  \item Endpoints of a maximal matching always form a vertex cover;
        endpoints of a maximum matching always form a minimum cover.
  \item The indicators of two edges sharing a vertex must be independent
        in order to sum their expectations.
\end{enumerate}

\clearpage
\subsection{Answer sketches}
\paragraph{Q1. Scheduling.}
Use layers
\[
s\longrightarrow i\longrightarrow(i,j)\longrightarrow d\longrightarrow t
\]
with capacities $b_i,1,1,1$, respectively.  Include $(i,j)\to d$ exactly
when task $d$ is in period $j$ and $d\in A_i$.  The target flow value is
the number $D$ of tasks.  A schedule sends one unit on the path for each
assignment.  Conversely, an integral flow of value $D$ saturates all
task-to-sink arcs.  Each task then has exactly one unit of incoming flow;
decode that arc as its assigned worker.  The worker and worker-period
capacities enforce the two quotas.  Fractional flow can split a task across
workers, but integer capacities guarantee the existence of an integral
maximum flow, which the standard augmenting-path algorithms can return.

For available pairs use binary $a_{id}$ and the feasibility ILP
\[
\sum_{i:d\in A_i}a_{id}=1\quad(\text{each }d),\qquad
\sum_{d\in A_i}a_{id}\le b_i\quad(\text{each }i),\qquad
\sum_{\substack{d\in A_i\\d\text{ in period }j}}a_{id}\le1
\quad(\text{each }i,j).
\]
The network has polynomial size and requires one maximum-flow computation.

\paragraph{Q2. Bounded-frequency covering.}
Minimize $\sum_j c_jx_j$ subject to
$\sum_{j:e\in S_j}x_j\ge1$ for every element, and $x_j\in\{0,1\}$.
Relax to $[0,1]$ and select $x_j^*\ge1/k$.  A covering constraint has at
most $k$ terms, so at least one meets the threshold.  With nonnegative costs,
\[
\ALG\le\sum_j kc_jx_j^*=kLP^*\le k\OPT.
\]
Bounded \emph{set size} does not bound the number of terms in a covering
constraint.  For example, one element and two identical unit-cost singleton
sets admit optimal LP solution $(1/2,1/2)$.  Set size is at most $1$, but
threshold $1$ selects neither set.

\paragraph{Q3. \texorpdfstring{MAX-$k$-CUT}{MAX-k-CUT} and conditional expectation.}
Independent uniform vertex labels cut an edge with probability $1-1/k$.
Linearity gives expectation $(1-1/k)\sum_e w_e\ge(1-1/k)\OPT$; edge
indicators need not be independent.  Under a partial labelling, an edge
with both endpoints fixed contributes its weight or zero.  An edge with
at least one endpoint unfixed contributes $(1-1/k)w_e$.  For the next
vertex, only edges to already fixed neighbors have a label-dependent
contribution.  Choose the label maximizing the weight of differently
labelled fixed neighbors.  The current conditional expectation is the
average of the $k$ branch expectations, so this choice cannot decrease it.
After all choices the objective equals its conditional expectation.
A direct implementation takes $O(k(n+m))$ time.

\paragraph{Q4. Set cover with rejection sampling.}
For any element $e$, one round misses it with probability
\[
\prod_{j:e\in S_j}(1-x_j^*)
\le \exp\!\left(-\sum_{j:e\in S_j}x_j^*\right)\le e^{-1}.
\]
After $r$ rounds this is at most $e^{-r}\le1/(2n)$.  A union bound implies
trial success probability $p\ge1/2$.  For trial cost $C$,
\[
\mathbb E[C]=\sum_j c_j\bigl(1-(1-x_j^*)^r\bigr)
\le r\sum_j c_jx_j^*=rLP^*.
\]
The accepted trial has the distribution of one trial conditioned on success
$A$.  Since $C\ge0$,
\[
\mathbb E[C\mid A]
=\frac{\mathbb E[C\mathbf1_A]}p
\le\frac{\mathbb E[C]}p\le2rLP^*\le2r\OPT.
\]
Fresh independent trials give a geometric stopping time of mean $1/p\le2$;
nontermination probability is $\lim_{t\to\infty}(1-p)^t=0$.
Each trial is polynomial time and its output is checked for feasibility.
The success event may be correlated with cost, so the unconditioned trial
expectation alone does not bound the returned cost.

\paragraph{Q5. Cuts.}
For a fixed maximum flow and any minimum-cut source side $S$,
\[
0=c(\delta^+(S))-|f|
=\sum_{e\in\delta^+(S)}(c_e-f_e)
 +\sum_{e\in\delta^-(S)}f_e.
\]
All summands are nonnegative, so each vanishes.  An arc with $c_e>0$
crossing outwards forces $f_e=c_e>0$; the same arc crossing inwards for
another minimum cut would force $f_e=0$, a contradiction.
For tie-breaking let $m=|E|$ and set $c'_e=(m+1)c_e+1$.  Then
\[
c'(\delta^+(S))=(m+1)c(\delta^+(S))+|\delta^+(S)|.
\]
An integer capacity difference of at least $1$ outweighs every cardinality
difference, which is at most $m$.  Minimize the new capacity with one
max-flow/min-cut computation.  Its encoding length grows by $O(\log(m+1))$.

\paragraph{Q6. MAX-DICUT.}
By independence for the two distinct endpoints and the LP constraints,
\begin{align*}
\Pr[i\in U,j\notin U]
&=(1/4+x_i/2)(1/4+(1-x_j)/2)\\
&\ge(1/4+z_{ij}/2)^2
=z_{ij}/2+(2z_{ij}-1)^2/16\ge z_{ij}/2.
\end{align*}
Sum with the nonnegative weights.  For direct rounding, put
$L=\sum w_{ij}z_{ij}=LP^*$ and $W=\sum w_{ij}$.  If $W=0$ the claim is
trivial.  Otherwise $x_i(1-x_j)\ge z_{ij}^2$, and Cauchy--Schwarz gives
\[
\mathbb E[\ALG]\ge\sum w_{ij}z_{ij}^2\ge L^2/W.
\]
The feasible point $x_i=z_{ij}=1/2$ implies $L\ge W/2$ when the LP
solution is optimal.  Hence $\mathbb E[\ALG]\ge L/2\ge\OPT/2$.
The shifted rule has a local per-arc proof; the direct rule uses LP
optimality for its global comparison.

\paragraph{Q7. MAX-SAT.}
Remove duplicate literals; treat tautological clauses as always satisfied.
For each remaining nonempty clause $C_i$, use
\[
\max\sum_iw_iy_i,\qquad
\sum_{j\in C_i^+}x_j+\sum_{j\in C_i^-}(1-x_j)\ge y_i,
\qquad x_j,y_i\in\{0,1\};
\]
relax the domains to $[0,1]$.  Independent variable rounding gives literal
success probabilities $a_1,\ldots,a_\ell$ in a clause with distinct
variables, and $\sum_h a_h\ge y_i^*$.  AM--GM gives
\[
\Pr[C_i\text{ fails}]=\prod_h(1-a_h)
\le\left(1-\frac{\sum_h a_h}{\ell}\right)^\ell
\le(1-y_i^*/\ell)^\ell.
\]
The function $g(y)=1-(1-y/\ell)^\ell$ is concave on $[0,1]$ with
$g(0)=0$.  Thus $g(y)\ge yg(1)$, and the clause contributes at least
$b_\ell w_i y_i^*$ where $b_\ell=1-(1-1/\ell)^\ell\ge1-1/e$.
Fair assignment contributes at least $a_\ell w_i y_i^*$, with
$a_\ell=1-2^{-\ell}$ because $y_i^*\le1$.

For $\ell=1,2$, $a_\ell+b_\ell=3/2$; for $\ell\ge3$,
$a_\ell\ge7/8$ and $b_\ell\ge1-1/e>5/8$.  Therefore, for outputs $A,B$,
\[
\mathbb E[\max(A,B)]\ge\tfrac12(\mathbb E[A]+\mathbb E[B])
\ge\tfrac34 LP^*\ge\tfrac34\OPT.
\]
The two outputs need not be independent.  Averaging just the global
$1/2$ and $1-1/e$ guarantees would not prove $3/4$.

\paragraph{Q8. TSP.}
Deleting a tour edge proves $c(T)\le\OPT$ for an MST $T$.  Doubling gives
even degrees and cost $2c(T)$; triangle inequality makes the final shortcut
tour no more expensive.  For Christofides, the odd-degree set has even
cardinality.  Shortcut an optimal tour to this set and split its alternating
edges into two perfect matchings.  The minimum matching $M$ costs no more
than their average, at most $\OPT/2$.  Thus the output costs at most
$c(T)+c(M)\le3\OPT/2$.

For the requested example, use the unit-cost cycle $A,B,D,C,A$ and diagonal
costs $2$.  MST $\{AB,AC,CD\}$ has Euler circuit $A,B,A,C,D,C,A$ after
doubling.  First-visit shortcutting returns $A,B,C,D,A$, cost $6$, while
the original cycle has optimal cost $4$.  Without triangle inequality a
shortcut edge can cost arbitrarily more than the detour it replaces.

\paragraph{Q9. Knapsack.}
Discard $w_i>W$, and handle no remaining items or $V=0$ separately.
Let $D[i,p]$ be the minimum weight of a subset of the first $i$ items with
value exactly $p$.  Set $D[0,0]=0$, all other initial states to $\infty$, and
\[
D[i,p]=\min\{D[i-1,p],\ w_i+D[i-1,p-v_i]\},
\]
omitting the second term when $p<v_i$.  Fill increasing $i$; choose the
largest $p$ with $D[n,p]\le W$ and backtrack stored choices.  Since
$p\le nV$, this takes $O(n^2V)$ arithmetic operations.

For $0<\epsilon<1$, put $K=\epsilon V/n$ and
$\widehat v_i=\lfloor v_i/K\rfloor\le n/\epsilon$.  Solve the DP with
scaled values and original weights.  It takes $O(n^3/\epsilon)$ operations.
For its returned set $S$ and an original optimum $S^*$,
\[
v(S)\ge K\widehat v(S)\ge K\widehat v(S^*)
\ge\OPT-nK=\OPT-\epsilon V\ge(1-\epsilon)\OPT.
\]
Filtering ensures that the maximum-value item is itself feasible, so
$V\le\OPT$.  Keeping original weights preserves feasibility.

\paragraph{Q10. Vertex form of Menger's theorem.}
Split each internal vertex into $v_{\rm in}\to v_{\rm out}$ of capacity
$1$, leaving $x,y$ unsplit.  Replace each undirected edge $\{u,v\}$ by
$u_{\rm out}\to v_{\rm in}$ and $v_{\rm out}\to u_{\rm in}$ of capacity
$M=n+1$, interpreting an endpoint $x$ or $y$ as the unsplit vertex.
Internally disjoint paths induce a feasible integral flow.  Conversely,
decompose integral flow into unit $x$--$y$ paths and cycles and discard
cycles; capacity-$1$ splitting arcs ensure internal vertex-disjointness.

A separator of size $k$ deletes $k$ splitting arcs and destroys all flow
paths.  Reachability after this deletion gives an original-network cut
of capacity at most $k$.  Because $x,y$ are nonadjacent, all $n-2$ internal
vertices form a separator, so a minimum cut has capacity at most
$n-2<M$.  Consequently it cuts no capacity-$M$ arc.  Its outgoing splitting
arcs identify a vertex separator whose size equals the cut capacity.
Max-flow/min-cut now equates the two optima.  Adjacency is excluded because
a direct $x$--$y$ edge survives deletion of every internal vertex.

\paragraph{Q11. Flow algorithms.}
A two-arc path with capacities $3,1$ has maximum flow $1$, but thresholds
$3,1.5$ find no path before the flawed algorithm stops.  Start instead at
$2^{\lfloor\log_2 C\rfloor}$ and halve through threshold $1$; handle $C=0$
separately.  Use only residual arcs of capacity at least the current threshold.

For BFS augmentation, let $d(v)$ be residual distance from $s$ before an
augmentation.  Every surviving old arc satisfies $d(v)\le d(u)+1$.
A new arc $(v,u)$ reverses an arc of the chosen shortest path, so
$d(v)=d(u)+1$; it too satisfies the inequality needed along any new path.
Induction along a new residual path proves that no distance decreases.
If $(u,v)$ is saturated on a shortest path, then $d(v)=d(u)+1$.
Before it can be saturated in that direction again, an augmentation must
use $(v,u)$; at that later time
\[
d_{\rm later}(u)=d_{\rm later}(v)+1
\ge d_{\rm earlier}(v)+1=d_{\rm earlier}(u)+2.
\]
Thus each labelled residual arc can be critical only $O(n)$ times.
Each augmentation saturates one such arc, giving $O(nm)$ augmentations
and $O(nm^2)$ time under the usual $O(m)$ search bound.

With integer capacities, arbitrary Ford--Fulkerson increases value by at
least $1$, giving $O(m|f^*|)$ time.  The numeric value $|f^*|$ may be
exponential in its binary encoding length.  With irrational capacities,
there need be no positive lower bound on augmentation amounts; bad path
choices can produce infinitely many positive augmentations.  BFS avoids
this dependence by bounding the number of augmentations combinatorially.

\paragraph{Q12. Common subsequence.}
For each symbol $a$, compute $f_a=\min_i\operatorname{count}_a(s_i)$.
Return $a$ repeated $f_a$ times for a symbol maximizing $f_a$; this is
always a common subsequence.  Some symbol appears at least $\OPT/q$ times
in an optimal common subsequence, and its $f_a$ is at least that large.
Alternatively,
\[
\OPT\le\sum_a f_a\le q\max_a f_a.
\]
To achieve $O(nL+q)$ even for variable $q$, initialize arrays once; while
scanning each string, use timestamps to count only its occurring symbols.
For each touched symbol update its minimum count and the number of strings
in which it appears.  Symbols appearing in fewer than $n$ strings get
$f_a=0$.  Each character and each touched symbol costs $O(1)$, followed
by an $O(q)$ final scan; no $O(q)$ reset per string is needed.

\subsection{True/false answers}
\begin{enumerate}[label=\textbf{T\arabic*.}]
  \item \textbf{False.} Take $q$ disjoint edges and one isolated vertex.
        Selecting all edge endpoints is a 2-approximate cover; its complement
        has size $1$, while a maximum independent set has size $q+1$.
  \item \textbf{False.} Use unit arcs $s\to a,s\to b,a\to c,b\to c,c\to t$.
        Sending $1/2$ on each branch and $1$ on $c\to t$ is a fractional
        maximum flow.  An integral maximum also exists.
  \item \textbf{True.} Weak duality gives $|f|\le|f^*|\le c(S)$;
        equality of the outer values certifies maximum flow and minimum cut.
  \item \textbf{False.} Capacity $C$ needs only $O(\log(C+1))$ bits.
        A bound polynomial in $C$ may be exponential in input length.
  \item \textbf{False.} A satisfiable formula can have minimum unsatisfied
        cost $0$.  An assignment satisfying only half of its clauses may
        meet a MAX-SAT factor-$2$ guarantee while leaving positive cost,
        impossible for any finite multiplicative guarantee relative to $0$.
  \item \textbf{False.} Use unit arcs along $s\to a\to b\to c\to d\to t$
        and an additional arc $s\to d$.  Initially $d(s,c)=3$.  Augment
        on the long path: the unused arc $s\to d$ and new reverse arc
        $d\to c$ give distance $2$ to $c$.
  \item \textbf{False.} The optimal relaxation value is a \emph{lower}
        bound.  An arbitrary fractional feasible value need not bound the
        integral optimum in either direction.  An integral feasible solution
        does give an upper bound.
  \item \textbf{False.} $K_n$ witnesses a worst-case gap $2-2/n$:
        integral optimum $n-1$, LP optimum $n/2$.  This is not a statement
        about every graph; a single edge with isolated vertices has gap $1$.
  \item \textbf{First true, second false.} An uncovered edge could be added
        to a matching, contradicting maximality.  For one edge the maximum
        matching's endpoints give cover size $2$, but minimum cover size $1$.
  \item \textbf{False.} Linearity of expectation requires no independence.
\end{enumerate}
